% Generated draft: requires mathematical and source review.
\documentclass[11pt]{article}
\usepackage[margin=25mm]{geometry}
\usepackage{fontspec}
\setmainfont{Latin Modern Roman}
\usepackage{amsmath,amssymb,mathtools,unicode-math}
\setmathfont{Latin Modern Math}
\usepackage{enumitem,needspace,xurl,hyperref}
\setlength{\parindent}{0pt}
\setlength{\parskip}{5pt}
\setlength{\emergencystretch}{3em}
\newcommand{\N}{\mathbb{N}}
\newcommand{\Z}{\mathbb{Z}}
\newcommand{\Q}{\mathbb{Q}}
\newcommand{\R}{\mathbb{R}}
\newcommand{\C}{\mathbb{C}}
\newcommand{\F}{\mathbb{F}}
\newcommand{\A}{\mathbb{A}}
\newcommand{\PP}{\mathbb{P}}
\newcommand{\E}{\mathbb{E}}
\newcommand{\eps}{\varepsilon}
\newcommand{\dd}{\,\mathrm{d}}
\begin{document}
\section*{UnsolvedMath ID: 6700022}
\textbf{Source problem code:} catalogue entry
\subsection{Definitions and mathematical statement}
The principal question is to evaluate the source's two invariants \(\sigma_\circ(X_0)\) and \(\sigma_\square(X_0)\) for simple Riemannian manifolds \(X_0=(X,g_0)\). Their definitions are not included in the record; the circular and square decorations are retained as different source symbols, not replaced by the Yamabe invariant or an invented extremal formula. Simple is also not a defined manifold class.

The source reports the equality
\[
\sigma_\circ(X_0)=\sigma_\square(X_0)=\operatorname{Sc}(X_0)
\]
for many compact symmetric spaces, and conjectures it for all intended compact symmetric spaces. Which cases and hypotheses are covered depends on the missing invariant definitions and source section 17. It asks about the contrast between Dirac-operator methods and minimal-hypersurface methods in obtaining sharp evaluations.

The concrete sphere statement quoted as known by Llarull's theorem is: for the unit round metric \(g_{\mathrm{rd}}\) on \(S^n\), if a smooth metric \(g\) satisfies
\[
\forall x\in S^n\ \forall v\in T_xS^n,\qquad
g_x(v,v)\ge(g_{\mathrm{rd}})_x(v,v),
\]
then
\[
\exists x\in S^n,\qquad \operatorname{Sc}_g(x)\le n(n-1).
\]
The open methodological request is for a proof by minimal hypersurfaces for \(n\ge3\), not a claim that the inequality itself is unknown. The source regards dimension three as plausible and complex/quaternionic projective-space analogues as less likely for that method. It also mentions incomplete examples with sharp minimal-hypersurface evaluations in section 21; this does not establish impossibility of all Dirac-based proofs.

The extraction then starts Schick's 1998 example: an \(n\)-torus modified by a two-handle/circle surgery on three times a fundamental-group generator, reported as a spin Schoen--Yau--Schick manifold for \(n\ge4\) with
\[
\pi_1(X_{\mathrm{Sch}})\cong\Z^{n-1}\times\Z/3\Z.
\]
The surgery/framing and named property are not defined, and the stored passage ends mid-sentence after all (known?. No missing conclusion about indices or scalar curvature is reconstructed. The supplied research summary labels the invariant evaluations and minimal-hypersurface comparison only partially resolved.
\subsection{Short English statement}
Evaluate the two source-specific scalar-curvature invariants and compare sharp Dirac and minimal-hypersurface methods. The sphere inequality is reported known; the requested hypersurface proof and broader evaluations remain qualified, and the Schick example is truncated.
\subsection{Sources}
\begin{itemize}
\item Local UnsolvedMath registry, permanent ID \(6700022\), exact source code \texttt{catalogue entry}; printed scalar-curvature tag [?21], occurrence 22. Proposer Misha Gromov, 2017.
\item Gromov, \emph{101 Questions, Problems and Conjectures around Scalar Curvature} (2017), recorded PDF page 30; \url{https://www.ihes.fr/~gromov/wp-content/uploads/2018/08/101-problemsOct1-2017.pdf}. Recorded access 2026-07-29; pdf-text-normalized extraction. Cross-references to sections 17 and 21 lack the referenced definitions/examples here.
\item Llarull's theorem is explicitly quoted in the source. The supplied research summary additionally names Goette and Gromov in connection with sharp Dirac bounds, without complete publication details.
\item Thomas Schick example (1998), source reference [105], whose bibliography entry and the end of the quoted passage are absent.
\end{itemize}
\end{document}
